\item Considere el oscilador amortiguado que se muestra en la figura 15.19. La masa del objeto es $375\text{ g}$, la constante de resorte es $100\text{ N/m}$ y $b = 0.100\text{ N}\cdot\text{s/m}$.
\begin{enumerate}
    \item ¿Durante qué intervalo de tiempo la amplitud disminuye a la mitad de su valor inicial?
    \item ¿Qué pasaría si? ¿Durante qué intervalo de tiempo la energía mecánica disminuye a la mitad de su valor inicial?
    \item Demuestre que, en general, la relación fraccionaria a la cual la amplitud disminuye en un oscilador armónico amortiguado es la mitad de la relación fraccionaria a la que disminuye la energía mecánica.
\end{enumerate}